\documentclass[10pt]{article}
\usepackage[T1]{fontenc}
\usepackage[utf8]{inputenc}
\usepackage[french]{babel}
\usepackage{amsmath,amssymb,mathrsfs,mathpazo}
\usepackage{geometry}
\geometry{a4paper,margin=2cm}
\usepackage{graphicx,array,tabularx,tabularray,colortbl,multirow,multicol,booktabs,enumitem,pifont,ulem,fancybox,diagbox,cancel}
\usepackage{pstricks,pstricks-add,pst-plot,pst-tree,pst-node,pst-text,pst-func,pst-math,pst-bspline,pst-3dplot,pst-eucl,pst-solides3d}
\let\pscalculate\undefined
\usepackage{pst-bezier}
\usepackage{tikz,pgfplots,tkz-tab}
\usetikzlibrary{arrows,arrows.meta,patterns,calc,positioning,trees,decorations.markings,shapes.geometric,shapes.misc,shapes.arrows,intersections,decorations.pathmorphing,decorations.pathreplacing,backgrounds,shadows,babel}
\pgfplotsset{compat=1.18}
\usepgfplotslibrary{fillbetween}
\usepackage[np]{numprint}
\usepackage{listings,makecell,pas-tableur,twoopt}
\let\barre\undefined
\usepackage[pstricks]{tablvar}
\definecolor{darkgreen}{RGB}{0,100,0}
\providecommand{\red}{\color{red}}
\providecommand{\blue}{\color{blue}}
\providecommand{\green}{\color{green}}
\providecommand{\black}{\color{black}}
\providecommand{\white}{\color{white}}
\providecommand{\gray}{\color{gray}}
\providecommand{\cyan}{\color{cyan}}
\providecommand{\magenta}{\color{magenta}}
\providecommand{\yellow}{\color{yellow}}
\providecommand{\violet}{\color{violet}}
\providecommand{\orange}{\color{orange}}
\providecommand{\euro}{\texteuro}
\providecommand{\EUR}{\texteuro}
\providecommand{\ieme}{\textsuperscript{e}}
\providecommand{\degres}{\ensuremath{^\circ}}
\providecommand{\degree}{\ensuremath{^\circ}}
\providecommand{\vv}[1]{\overrightarrow{#1}}
\providecommand{\Coord}[1]{#1}
\providecommand{\decosix}{}
\providecommand{\starredbullet}{\ensuremath{\bullet}}
\pagestyle{empty}
\setlength{\parindent}{0pt}
\usepackage{adjustbox}\DeclareUnicodeCharacter{2236}{:}\DeclareUnicodeCharacter{2190}{\ensuremath{\leftarrow}}
\providecommand{\R}{\mathbb{R}}
\providecommand{\N}{\mathbb{N}}
\providecommand{\Z}{\mathbb{Z}}
\providecommand{\Q}{\mathbb{Q}}
\providecommand{\D}{\mathbb{D}}
\providecommand{\C}{\mathbb{C}}
\providecommand{\e}{{\mathrm{e}}}
\providecommand{\vect}[1]{\overrightarrow{#1}}
\providecommand{\vectt}[1]{\overrightarrow{\mathrm{#1}}}
\providecommand{\vvt}[1]{\overrightarrow{\mathrm{#1}}}
\providecommand{\np}[2][]{{#2~#1}}
\providecommand{\vv}[1]{{\overrightarrow{#1}}}
\renewcommand{\d}{{\mathrm{d}}}
\providecommand{\degres}{{^\circ}}
\providecommand{\degre}{{^\circ}}
\providecommand{\blue}{\color{blue}}
\providecommand{\red}{\color{red}}
\providecommand{\enskip}{\quad}
\providecommand{\llbracket}{\lbrack\!\lbrack}
\providecommand{\rrbracket}{\rbrack\!\rbrack}
\providecommand{\euro}{{\text{€}}}
\providecommand{\no}{{\text{n°}}}
\providecommand{\interval}[2]{{[#1;#2]}}
\providecommand{\intervFO}[2]{{[#1;#2[}}
\providecommand{\dotfill}{\ldots}

\providecommand{\ds}{\displaystyle}
\providecommand{\textding}[1]{#1}
\providecommand{\Oij}{\(\mathrm{O};\vec{\imath},\vec{\jmath}\)}
\providecommand{\Oijk}{\(\mathrm{O};\vec{\imath},\vec{\jmath},\vec{k}\)}
\providecommand{\Ouv}{\(\mathrm{O};\vec{u},\vec{v}\)}
\providecommand{\intervFF}[2]{\ensuremath{[#1;#2]}}
\renewcommand{\euro}{\ifmmode\text{\texteuro}\else\texteuro\fi}
\DeclareUnicodeCharacter{00B0}{\ifmmode^{\circ}\else\textdegree\fi}
\DeclareUnicodeCharacter{20AC}{\ifmmode\text{\texteuro}\else\texteuro\fi}
% Énoncés : Éducation nationale. Transcription et corrigés : APMEP, Denis Vergès et contributeurs. Découpage et classement : CoopMaths, Sébastien Lozano, Éric Elter et contributeurs. Adaptation de la présentation : Maths Explorer.
% CC BY-SA : https://coopmaths.fr/www/about/
% https://coopmaths.fr/alea/static/eam/2025/tex/eam_specialite_2025_06_sujet0v1_1_cor.tex
\begin{document}
\noindent\textit{APMEP / CoopMaths. CC BY-SA. Présentation adaptée par Maths Explorer.}\par\medskip
\medskip

\textbf{Question 1}

\medskip

L'inverse du double de 5 est égal à :

\begin{center}
\begin{tabularx}{\linewidth}{*{4}{X}}
\textbf{a.~} $\dfrac 25$ &\textbf{b.~} $\dfrac{1}{10}$&\textbf{c.~} $\dfrac 52$ &\textbf{d.~} 10
\end{tabularx}
\end{center}

\begin{tabular}{@{\hspace{0.03\linewidth}} | p{0.95\linewidth}}
Le double de 5 est 10, l'inverse du double de 5 est donc $\dfrac{1}{10}$.
\hfill\textbf{Réponse b}
\end{tabular}

\bigskip

\textbf{Question 2}

\medskip

On considère la relation $F = a + \dfrac{b}{cd}.$\\[5pt]
Lorsque $a = \dfrac12$, $b = 3$, $c = 4$, $d = - \dfrac14$, la valeur de $F$ est égale à :

\begin{center}
\begin{tabularx}{\linewidth}{*{4}{X}}
\textbf{a.~} $- \dfrac52$&\textbf{b.~}$- \dfrac32$&\textbf{c.~} $\dfrac52$&\textbf{d.~} $\dfrac32$
\end{tabularx}
\end{center}

\begin{tabular}{@{\hspace{0.03\linewidth}} | p{0.95\linewidth}}
$cd=4\times \left (-\dfrac{1}{4}\right )=-1$ donc $\dfrac{b}{cd}=\dfrac{3}{-1}=-3$

$F= a+\dfrac{b}{cd} = \dfrac{1}{2} + (-3)=\dfrac{1}{2}-\dfrac{6}{2}=-\dfrac{5}{2}$
\hfill\textbf{Réponse a}
\end{tabular}

\bigskip

\textbf{Question 3}

\medskip

Le prix d'un article est multiplié par 0,975.

Cela signifie que le prix de cet article a connu :

\begin{center}
\begin{tabularx}{\linewidth}{*{2}{X}}
\textbf{a.~} une baisse de 2,5\,\% &\textbf{b.~} une augmentation de 97,5\,\%\\
\textbf{c.~} une baisse de 25\,\%& \textbf{d.~} une augmentation de 0,975\,\%
\end{tabularx}
\end{center}

\begin{tabular}{@{\hspace{0.03\linewidth}} | p{0.95\linewidth}}
Multiplier par $\left (1-\dfrac{t}{100}\right )$, c'est faire subir une baisse de $t\,\%$.

$0,975 = 1-0,025 = 1 - \dfrac{2,5}{100}$; il s'agit donc d'une baisse de $2,5\,\%$.

\hfill\textbf{Réponse a}
\end{tabular}

\bigskip 

\textbf{Question 4}

\medskip

Le prix d'un article est noté $P$. Ce prix augmente de 10\,\% puis baisse de 10\,\%.

À l'issue de ces deux variations, le nouveau prix est noté $P_1$. On peut affirmer que:

\begin{center}

\textbf{a.~} $P_1 = P$\hfill \textbf{b.~} $P_1 > P$\hfill
\textbf{c.~} $P_1 < P$\hfill \textbf{d.~} Cela dépend de $P$

\end{center}

\begin{tabular}{@{\hspace{0.03\linewidth}} | p{0.95\linewidth}}
Augmenter de 10\,\%, c'est multiplier par $1+\dfrac{10}{100}=1,1$.

Baisser de 10\,\%, c'est multiplier par $1-\dfrac{10}{100}=0,9$.

Augmenter de 10\,\% puis baisser de 10\,\%, c'est multiplier par $1,1\times 0,9 = 0,99<1$.

Donc $P_1<P$\hfill\textbf{Réponse c}
\end{tabular}

\bigskip 

\textbf{Question 5}

\medskip

On lance un dé à 4 faces. \\
La probabilité d'obtenir chacune des faces est donnée dans le tableau ci-dessous :

\begin{center}
\begin{tabularx}{\linewidth}{|*{4}{>{\centering \arraybackslash}X|}}
\hline
Face numéro 1 &Face numéro 2& Face numéro 3& Face numéro 4\\ 
\hline
 0,5 &$\dfrac16$&0,2& $x$\rule[-10pt]{0pt}{28pt}\\
 \hline
\end{tabularx}
\end{center}

\medskip

On peut affirmer que :

\begin{center}
\begin{tabularx}{\linewidth}{*{4}{X}}
\textbf{a.~} $x = \dfrac{2}{15}$&\textbf{b.~} $x = \dfrac{2}{3}$& 
\textbf{c.~} $x = 0,4$&\textbf{d.~} $x = 0,1$
\end{tabularx}
\end{center}

\begin{tabular}{@{\hspace{0.03\linewidth}} | p{0.95\linewidth}}
La somme des probabilités doit être égale à 1, donc
$0,5+\dfrac{1}{6}+ 0,2+x=1$

donc
$0,7+\dfrac{1}{6}+x=1$
donc
$\dfrac{7}{10}+\dfrac{1}{6}+x=1$
donc
$\dfrac{21}{30}+ \dfrac{5}{30}+x=1$
donc
$\dfrac{26}{30}+x=1$

donc
$x= \dfrac{30}{30}-\dfrac{26}{30}$
donc
$x=\dfrac{4}{30}$
donc
$x=\dfrac{2}{15}$
\hfill\textbf{Réponse a}
\end{tabular}

\bigskip 

\textbf{Question 6}

\medskip

On considère $x$, $y$, $u$ des réels non nuls tels que  $\dfrac 1x + \dfrac 1y = \dfrac 1u$.

On peut affirmer que :

\begin{center}
\begin{tabularx}{\linewidth}{*{4}{X}}
\textbf{a.~} $u = \dfrac{xy}{x + y}$&\textbf{b.~} $u = \dfrac{x + y}{xy}$ &\textbf{c.~} $u = xy$ &\textbf{d.~} $u = x + y$
\end{tabularx}
\end{center}

\begin{tabular}{@{\hspace{0.03\linewidth}} | p{0.95\linewidth}}
$\dfrac{1}{x}+\dfrac{1}{y}= \dfrac{1}{u}$
équivaut à
$\dfrac{x+y}{xy}=\dfrac{1}{u}$
équivaut à
$\dfrac{xy}{x+y}=u$
\hfill\textbf{Réponse a}
\end{tabular}

\bigskip

\begin{minipage}{0.65\linewidth}
\textbf{Question 7}

\medskip

On a représenté ci-contre la parabole d'équation $y = x^2$.

On note $(\mathcal{J})$ l'inéquation, sur $\R$,\: $x^2 \geqslant 10$.

L'inéquation $(\mathcal{J})$ est équivalente à :

{\renewcommand{\arraystretch}{1.5}
\begin{tabularx}{\linewidth}{*{24}{X}}
\textbf{a.~}$- \sqrt{10} \leqslant x \leqslant \sqrt{10}$&\textbf{b.~}$x \leqslant- \sqrt{10}$ ou $x \geqslant \sqrt{10}$\\
\textbf{c.~}$x \geqslant \sqrt{10}$&\textbf{d.~}$x = \sqrt{10}$ ou $x = - \sqrt{10}$
\end{tabularx}}
\end{minipage}\hfill
\begin{minipage}{0.32\linewidth}
\psset{xunit=0.5cm,yunit=0.2cm,arrowsize=2pt 3}
\begin{pspicture}(-5,-3)(5,16)
\psaxes[linewidth=1.25pt,Dx=20,Dy=10,labelFontSize=\scriptstyle]{->}(0,0)(-5,-1)(5,16)
\uput[u](4.8,0){$x$}\uput[r](0,15){$y$}
\psplot[plotpoints=1000,linewidth=1pt,linecolor=red]{-4}{4}{x dup mul}
\psplot[plotpoints=1000,linewidth=2pt,linecolor=blue]{-4}{-3.162}{x dup mul}
\psplot[plotpoints=1000,linewidth=2pt,linecolor=blue]{3.162}{4}{x dup mul}

\psset{linecolor=blue}
\psline[linewidth=2pt]{-]}(-5,0)(-3.162,0)
\psline[linewidth=2pt]{[-}(3.162,0)(5,0)
\psset{linestyle=dashed}
\psline(-5,10)(5,10)
\psline(-3.162,10)(-3.162,0)\uput{8pt}[d](-3.162,0){\blue\tiny$-\sqrt{10}$}
\psline(3.162,10)(3.162,0) \uput{8pt}[d](3.162,0){\blue\tiny$\sqrt{10}$}
\end{pspicture}
\end{minipage}

\medskip

\begin{tabular}{@{\hspace{0.03\linewidth}} | p{0.95\linewidth}}
$f(x)\geqslant 10$ signifie que la courbe représentant la fonction $f$ est au dessus de la droite d'équation $y=10$.

Voir graphique \hfill\textbf{Réponse b}
\end{tabular}


\textbf{Question 8}

\medskip

On a représenté ci-contre une droite $\mathcal{D}$ dans un repère orthonormé.

Une équation de la droite $\mathcal{D}$ est: 

\begin{minipage}{0.64\linewidth}
{\renewcommand{\arraystretch}{2}
\begin{tabularx}{\linewidth}{*{2}{X}}
\textbf{a.~} $y = - \dfrac32 x + 2$&\textbf{b.~} $y = \dfrac23 x + 2$\\
\textbf{c.~} $2x - 3y - 6 = 0$&\textbf{d.~} $\dfrac x3 + \dfrac y2 - 1 = 0$
\end{tabularx}}
\end{minipage}
\hfill
\begin{minipage}{0.32\linewidth}
\psset{unit=0.5cm,arrowsize=2pt 3}
\begin{pspicture*}(-3,-2.5)(6.5,4)
\psgrid[subgriddiv=1,  gridlabels=0, gridcolor=lightgray] 
\psaxes[linewidth=1.25pt,labelFontSize=\scriptstyle]{->}(0,0)(-2.99,-2.5)(6.5,4)
\psplot[plotpoints=500,linewidth=1.25pt,linecolor=blue]{-3}{6}{2 x 2 mul 3 div sub}
\uput[u](6,0){$x$}\uput[r](0,3.5){$y$}
\uput[dl](-1.5,3){\blue $\mathcal{D}$}
\end{pspicture*}
\end{minipage}

\medskip

\begin{tabular}{@{\hspace{0.03\linewidth}} | p{0.95\linewidth}}
Il suffit de chercher quelle droite passe par les points de coordonnées $(0\,,\,2)$ et $(3\,,\,0)$.

\hfill\textbf{Réponse d}
\end{tabular}

\bigskip

\textbf{Question 9}

\medskip

On considère trois fonctions définies sur $\R$ :

\begin{center}
\begin{tabularx}{\linewidth}{*{3}{X}}
$f_1 : \:x \longmapsto x^2 - (1 - x)^2$&$f_2 : x \longmapsto \dfrac x2 - \left(1 + \dfrac{1}{\sqrt{2}}\right)$&$f_3 : x \longmapsto \dfrac{5 - \frac23 x}{0,7}$
\end{tabularx}
\end{center}

Parmi ces trois fonctions, celles qui sont des fonctions affines sont:

\begin{center}
\begin{tabularx}{\linewidth}{*{2}{X}}
\textbf{a.~} aucune &\textbf{b.~} toutes\\
\textbf{c.~} uniquement la fonction $f_1$ &\textbf{d.~} uniquement les fonction $f_2$ et $f_3$
\end{tabularx}
\end{center}

\medskip

\begin{tabular}{@{\hspace{0.03\linewidth}} | p{0.95\linewidth}}
\begin{list}{\textbullet}{On cherche si les fonctions sont de la forme $x \longmapsto ax+b$.}
\item $f_1(x)=x^2-\left (1-x\right )^2=x^2 -\left (1-2x+x^2\right )=x^2-1+2x-x^2 = 2x-1$.

Donc $f_1$ est une fonction affine.
\item $f_2(x)= \dfrac{x}{2} - \left(1 + \dfrac{1}{\sqrt{2}}\right)$ est de la forme $ax+b$ avec $a=\dfrac{1}{2}$ et $b=-\left(1 + \dfrac{1}{\sqrt{2}}\right)$. 

Donc $f_2$ est une fonction affine.
\item $f_3(x)= \dfrac{5 - \frac23 x}{0,7} = -\dfrac{2}{3\times 0,7} x + \dfrac{5}{0,7}$ est de la forme $ax+b$ avec $a=-\dfrac{2}{2,1}$ et $b=\dfrac{5}{0,7}$. 

Donc $f_3$ est une fonction affine.
\end{list}

\hfill\textbf{Réponse b}
\end{tabular}

\bigskip

\textbf{Question 10}

\medskip

\begin{minipage}{0.65\linewidth}
On a représenté ci-contre une parabole $\mathcal{P}$.

Une seule des quatre fonctions ci-dessous est susceptible d'être représentée par la parabole $\mathcal{P}$.

Laquelle ?

\begin{center}
\begin{tabularx}{\linewidth}{*{2}{X}}
\textbf{a.~} $x \longmapsto x^2 - 10$ &\textbf{b.~} $x\longmapsto -x^2 - 10$ \\
\textbf{c.~} $x\longmapsto -x^2 + 10$ &\textbf{d.~} $x\longmapsto -x^2 + 10x$
\end{tabularx}
\end{center}
\end{minipage}\hfill
\begin{minipage}{0.33\linewidth}
\psset{xunit=0.6cm,yunit=0.2cm,arrowsize=2pt 3}
\begin{pspicture}(-3.5,-2)(3.5,12)
\psaxes[linewidth=1.25pt,Dx=10,Dy=15]{-}(0,0)(-3.5,-2)(3.5,12)
\psplot[plotpoints=1000,linewidth=1.25pt,linecolor=red]{-3.5}{3.5}{10 x dup mul sub}
\uput[dl](0,0){O}
\uput[r](2.5,6.25){\red $\mathcal{P}$}
\end{pspicture}
\end{minipage}

\begin{tabular}{@{\hspace{0.03\linewidth}} | p{0.95\linewidth}}
Le point de la courbe d'abscisse 0 a une ordonnée strictement positive.
\hfill\textbf{Réponse c}
\end{tabular}


\textbf{Question 11}

\medskip

\begin{minipage}{0.65\linewidth}
On a représenté ci-contre la courbe $\mathcal{C}$ d'une fonction $f$.

Les points A, B, R et S appartiennent à la courbe $\mathcal{C}$.

Leurs abscisses sont notées respectivement $x_{\text{A}},\: x_{\text{B}},\: x_{\text{R}}$
 et $x_{\text{S}}$.

L'inéquation $x \times f(x) > 0$ est vérifiée par : 

\medskip

\begin{tabularx}{\linewidth}{*{2}{X}}
\textbf{a.~} $x_{\text{A}}$ et $x_{\text{B}}$ &\textbf{b.~} $x_{\text{A}}$ et $x_{\text{R}}$ \\
\textbf{c.~} $x_{\text{A}}$ et $x_{\text{S}}$ &\textbf{d.~} $x_{\text{A}}, \: x_{\text{B}}$ et $x_{\text{S}}$
\end{tabularx}
\end{minipage}\hfill
\begin{minipage}{0.33\linewidth}
\psset{xunit=0.6cm,yunit=0.5cm,arrowsize=2pt 3}
\begin{pspicture*}(-3.5,-4.5)(3.5,3)
\psaxes[linewidth=1.25pt,Dx=10,Dy=15]{-}(0,0)(-3.5,-4.5)(3.5,3)
\psplot[plotpoints=1000,linewidth=1.25pt,linecolor=blue]{-3.5}{3.5}{x 3 exp x dup mul 2 mul sub 5 x mul sub 6 add 6 div}
\uput[dl](0,0){O}\uput[r](-3,-4){A}\uput[ul](-1.44,0.97){B}\uput[ur](0.4,0.624){R}\uput[dr](2.4,-0.616){S}
\uput[r](2.4,0.8){\blue $\mathcal{C}$}
\psdots(-3,-4)(-1.44,0.97)(0.4,0.624)(2.4,-0.616)
\end{pspicture*}
\end{minipage}

\bigskip

\begin{tabular}{@{\hspace{0.03\linewidth}} | p{0.95\linewidth}}
\begin{list}{\textbullet}{On calcule $x\times f(x)$ pour les 4 points.}
\item $x_{\text{A}} <0$ et $f\left (x_{\text{A}}\right )<0$ donc $x_{\text{A}} \times f\left (x_{\text{A}}\right )>0$, donc $x_{\text{A}}$ est solution.
\item $x_{\text{B}} <0$ et $f\left (x_{\text{B}}\right )>0$ donc $x_{\text{B}} \times f\left (x_{\text{B}}\right )<0$, donc $x_{\text{B}}$ n'est pas solution.
\item $x_{\text{R}} >0$ et $f\left (x_{\text{R}}\right )>0$ donc $x_{\text{R}} \times f\left (x_{\text{R}}\right )>0$, donc $x_{\text{R}}$ est solution.
\item $x_{\text{S}} >0$ et $f\left (x_{\text{S}}\right )<0$ donc $x_{\text{S}} \times f\left (x_{\text{S}}\right )<0$, donc $x_{\text{S}}$ n'est pas solution.
\end{list}

\hfill\textbf{Réponse b}
\end{tabular}

\bigskip

\textbf{Question 12}

\medskip

Voici une série de notes avec les coefficients associés.

\begin{center}
\begin{tabularx}{0.6\linewidth}{|l |*{3}{>{\centering \arraybackslash}X|}}\hline
Note&10 &8 &16\\ \hline
Coefficient &1 &2&$x$\\ \hline
\end{tabularx}
\end{center}

On note $m$ la moyenne de cette série.
Que doit valoir $x$ pour que $m = 15$ ? 

\begin{center}
\begin{tabularx}{\linewidth}{*{4}{X}}
\textbf{a.~} impossible &\textbf{b.~} $x = 10^{-3}$&
\textbf{c.~} $x = 3$ &\textbf{d.~} $x = 19$
\end{tabularx}
\end{center}

\begin{tabular}{@{\hspace{0.03\linewidth}} | p{0.95\linewidth}}
Il y a en tout $(x+3)$ notes, et si leur moyenne est de 15, la somme des notes doit être égale à $15\times (x+3)$ soit $15x+45$.

Or la somme des notes est égale à $10\times 1+8\times 2+16\times x = 16x+26$.

$16x+26 = 15x+45$ équivaut à $16x-15x=45-26$ soit $x=19$.
\hfill\textbf{Réponse d}
\end{tabular}

\end{document}
